\section{Application to \texttt{onto-med-rag}}

This section connects the concepts above to our competition project
(\textbf{Bài 2 --- Ontological Reasoning in Medical Knowledge Retrieval}):
Vietnamese clinical text in, structured entities with ICD-10/RxNorm codes out.

\subsection{Project goal in ontology terms}

The competition asks us to turn unstructured clinical notes into a \emph{knowledge
graph fragment} per document:

\begin{itemize}
  \item \textbf{Entities} --- symptoms, diagnoses, drugs, lab tests, results;
  \item \textbf{Relationships} --- (Phase 2+) links between concepts in the same note;
  \item \textbf{Shared meaning} --- Vietnamese surface forms mapped to ICD-10 and
        RxNorm via global identifiers;
  \item \textbf{Reasoning} --- assertions (negation, family, history) and, later,
        inferred links between concepts.
\end{itemize}

These are exactly the four patterns from Table~\ref{tab:kg-patterns}, applied to
clinical NLP rather than enterprise CRM/ERP data.

\subsection{Pipeline mapped to semantic layers}

Our Phase~1 pipeline (\texttt{docs/architecture.md}) and where ontology technology
fits:

\begin{table}[htbp]
  \centering
  \small
  \caption{Pipeline stages and ontology role}
  \label{tab:pipeline-ontology}
  \renewcommand{\arraystretch}{1.2}
  \begin{tabularx}{\textwidth}{@{}l l X@{}}
    \toprule
    \textbf{Stage} & \textbf{Module} & \textbf{Ontology connection} \\
    \midrule
    NER & \texttt{src/ner/} &
    Extract entity \emph{spans} --- the raw mentions that will become RDF subjects \\
    Classification & \texttt{src/classification/} &
    Assign RDF \emph{types} (\texttt{:Diagnosis}, \texttt{:Drug}, \texttt{:Symptom}, \ldots) \\
    Assertion detection & \texttt{src/assertion/} &
    Context flags as \emph{metadata on triples} (negated vs.\ active condition) \\
    Entity linking & \texttt{src/linking/} &
    Map Vietnamese text to ICD-10/RxNorm IRIs (\texttt{owl:sameAs}, synonyms) \\
    Candidate ranking & \texttt{src/ranking/} &
    Filter/rank using KB hierarchies (\texttt{rdfs:subClassOf}, disjointness) \\
    Output & \texttt{src/schemas/} &
    JSON entities $\approx$ a serialised ABox view of the graph \\
    \bottomrule
  \end{tabularx}
\end{table}

Phase~1 deliberately keeps ontology \emph{light}: ICD-10 and RxNorm act as structured
knowledge sources behind linking and ranking, not a full OWL reasoner
(\texttt{docs/architecture.md}). Phase~2+ adds relationship inference and deeper
ontological reasoning --- the competition name's long-term direction.

\newpage
\subsection{Output JSON as RDF triples}

Each predicted entity in \texttt{output/\{id\}.json} can be read as one or more
RDF triples. Take a diagnosis from our sample data
(\texttt{data/examples/sample\_output.json}):

\begin{verbatim}
# Surface form extracted by NER
:mention-trào-ngược a :Diagnosis ;
    :surfaceForm "bệnh trào ngược dạ dày - thực quản" ;
    :charOffset "99-134" .

# Linked to standard codes (candidates)
:mention-trào-ngược owl:sameAs icd:K21.0 .
:mention-trào-ngược owl:sameAs icd:K21.9 .

# ICD-10 hierarchy (from KB, not in JSON today)
icd:K21.0 rdfs:subClassOf icd:K21 .
icd:K21   rdfs:subClassOf icd:K20-K31 .
\end{verbatim}

A drug with historical assertion:

\begin{verbatim}
:mention-chlor a :Drug ;
    :surfaceForm "Chlorpheniramine 0.4 MG/ML" ;
    :hasAssertion :isHistorical .

:mention-chlor owl:sameAs rxnorm:360047 .
\end{verbatim}

The JSON schema (\texttt{src/schemas/entity.py}) is therefore a \textbf{projection}
of an RDF graph: fields \texttt{text}, \texttt{position}, \texttt{type},
\texttt{assertions}, and \texttt{candidates} map to literals, classes, annotation
properties, and external IRIs respectively.

\subsection{Entity types and knowledge bases}

\begin{table}[htbp]
  \centering
  \small
  \caption{Competition entity types mapped to ontology and KB}
  \label{tab:entity-kb}
  \renewcommand{\arraystretch}{1.15}
  \begin{tabularx}{\textwidth}{@{}l l l@{}}
    \toprule
    \textbf{JSON \texttt{type}} & \textbf{OWL class (example)} & \textbf{KB / codes} \\
    \midrule
    \texttt{CHẨN\_ĐOÁN} & \texttt{:Diagnosis} & ICD-10 (\texttt{data/kb/icd10/}) \\
    \texttt{THUỐC} & \texttt{:Drug} & RxNorm (\texttt{data/kb/rxnorm/}) \\
    \texttt{TRIỆU\_CHỨNG} & \texttt{:Symptom} & --- (assertions only) \\
    \texttt{TÊN\_XÉT\_NGHIỆM} & \texttt{:LabTest} & --- \\
    \texttt{KẾT\_QUẢ\_XÉT\_NGHIỆM} & \texttt{:LabResult} & --- \\
    \bottomrule
  \end{tabularx}
\end{table}

\textbf{Pattern 3 (shared meaning)} applies directly: ``trào ngược dạ dày - thực
quản'' (Vietnamese note) and ICD code \texttt{K21.0} (Gastro-esophageal reflux
disease with esophagitis) must refer to the same concept. RDF IRIs and
\texttt{owl:sameAs} are the standard way to unify surface text with standard
terminologies --- the same problem Bryon Jacob describes for \texttt{Client =
Customer = Account}, but with clinical vocabularies
(\href{https://bryon.io/why-rdf-is-the-natural-knowledge-layer-for-ai-systems-a5fd0b43d4c5}{Part~1 of the series}).

\newpage
\subsection{Assertions as semantic context}

Assertions change \emph{meaning}, not just labels:

\begin{table}[htbp]
  \centering
  \small
  \caption{Assertion flags and clinical interpretation}
  \label{tab:assertions}
  \renewcommand{\arraystretch}{1.15}
  \begin{tabularx}{\textwidth}{@{}l X@{}}
    \toprule
    \textbf{Assertion} & \textbf{Why an ontology layer helps} \\
    \midrule
    \texttt{isNegated} &
    ``không ho'' --- symptom present in text but \emph{not} attributed to patient;
    reasoner/QA must not treat as active finding \\
    \texttt{isFamily} &
    Condition belongs to a relative, not the patient --- disjoint patient vs.\ family
    context \\
    \texttt{isHistorical} &
    Past drug use or diagnosis --- temporal/context metadata on the triple \\
    \bottomrule
  \end{tabularx}
\end{table}

OWL can express constraints (e.g.\ a negated diagnosis is not an active
\texttt{:Diagnosis} for treatment planning). Even in Phase~1, encoding assertions
explicitly avoids the LLM ``guessing'' foreign-key-style semantics that the Bryon
series warns about for SQL schemas.

\subsection{Where RDFS and OWL help next}

\begin{table}[htbp]
  \centering
  \small
  \caption{Near-term ontology use by project phase}
  \label{tab:phase-roadmap}
  \renewcommand{\arraystretch}{1.15}
  \begin{tabularx}{\textwidth}{@{}l X@{}}
    \toprule
    \textbf{Phase} & \textbf{Ontology / tool use} \\
    \midrule
    Phase~1 (now) &
    Load ICD-10/RxNorm as lookup graphs; dictionary + retrieval linking;
    optional Turtle export of predictions for debugging \\
    Phase~2 &
    Infer relationships between entities in one note
    (\texttt{symptom} $\rightarrow$ \texttt{associatedWith} $\rightarrow$
    \texttt{diagnosis}); SPARQL over document graphs \\
    Phase~3+ &
    Owlready2 overlay ontology for Vietnamese clinical terms;
    \texttt{sync\_reasoner()} validation; LLM orchestrates NER while reasoner
    handles transitivity (e.g.\ drug class $\rightarrow$ ingredient) and consistency \\
    \bottomrule
  \end{tabularx}
\end{table}

Concrete OWL patterns from Section~3 that apply to medical linking:

\begin{itemize}
  \item \texttt{rdfs:subClassOf} --- ICD-10 hierarchy for ranking candidates
        (prefer specific codes, expand to parents when ambiguous);
  \item \texttt{owl:equivalentClass} / \texttt{owl:sameAs} --- Vietnamese synonyms
        and RxNorm concepts;
  \item \texttt{owl:TransitiveProperty} --- ``ingredient $\rightarrow$ drug class
        $\rightarrow$ therapeutic category'' in RxNorm;
  \item \texttt{owl:disjointWith} --- a \texttt{:Symptom} is not a \texttt{:Diagnosis}
        (classification sanity checks).
\end{itemize}

\newpage
\subsection{Owlready2 in our workflow}

For \texttt{onto-med-rag}, Owlready2 (Section~4) is the Python library for
anything we build beyond flat KB files:

\begin{enumerate}
  \item \texttt{uv add owlready2} --- add to \texttt{pyproject.toml}.
  \item Load or create \texttt{data/kb/clinical\_overlay.owl}; mirror ICD-10/RxNorm
        subsets relevant to the test set.
  \item Define classes for our five entity types and object properties for future
        relationship extraction; set Vietnamese \texttt{label}s on individuals.
  \item Map surface forms to ICD/RxNorm via \texttt{mapsToICD} /
        \texttt{mapsToRxNorm} properties or \texttt{owl:sameAs} equivalents.
  \item \texttt{sync\_reasoner()} --- catch unsatisfiable classes before linking
        code depends on the ontology.
  \item \texttt{onto.save()} and/or SPARQL queries from \texttt{src/linking/} and
        \texttt{src/ranking/}.
\end{enumerate}

Example sketch aligned with our JSON schema:

\begin{small}
\begin{verbatim}
from owlready2 import *

onto = get_ontology("http://onto-med-rag.example/clinical.owl")

with onto:
    class Diagnosis(Thing): pass
    class Drug(Thing): pass
    class icd10_code(Thing): pass
    class rxnorm_code(Thing): pass
    class hasCandidate(ObjectProperty): pass

    d = Diagnosis("mention_trào_ngược")
    d.label = ["bệnh trào ngược dạ dày - thực quản"]
    k21 = icd10_code("K21.0")
    d.hasCandidate = [k21]

onto.save("data/kb/clinical_overlay.owl")
\end{verbatim}
\end{small}

\subsection{End-to-end picture}

\begin{small}
\begin{verbatim}
Vietnamese clinical text (.txt)
        |
        v
   [NER + classifier]  -->  mentions with types (RDF subjects)
        |
        v
   [Assertion module]  -->  context on triples
        |
        v
   [Linking + ranking] -->  ICD-10 / RxNorm IRIs (global identity)
        |
        v
   output/{id}.json    -->  serialised ABox (Phase 1 deliverable)
        |
        v  (Phase 2+)
   RDF graph + OWL reasoner -->  inferred relationships, validated QA/RAG
\end{verbatim}
\end{small}

\textbf{Summary:} Sections~1--4 describe general patterns and tools; this project
instantiates them on Vietnamese EHR-style text, ICD-10, and RxNorm. Phase~1 ships
accurate JSON linking; the ontology stack is the path to Phase~2 relationship
inference and trustworthy medical retrieval without rebuilding RDF piecemeal.
